\section{本章小结}
\label{sec:besov-summary}

本章开头指出了一阶Sobolev条件的两处空白：它排除了阶跃函数，却读不出阶跃函数平移差的规律；它也没有规定不同尺度上的变化怎样汇总。平移差回答了第一点：周期阶跃在$L^2$中满足$\lVert\Delta_hf\rVert_2=(2|h|)^{1/2}$，光滑性可以是步长的分数次幂。同一尺度上，积分指数$p$决定位置不同的变化怎样度量，集中在小集合上的变化在大$p$下显得更剧烈；跨尺度上，统一上界看不出变化在多少个尺度上持续出现，需要单独的汇总方式。高阶差分抵消局部低阶多项式，连续模取步长不超过给定尺度的最坏情形；饱和现象说明$r$阶连续模至多表达$r$阶规律，所以刻画$s$阶光滑性须用$r>s$阶差分。

Besov范数把$t^{-s}\omega_r(f,t)_p$按$\mathrm dt/t$的$L^q$方式汇总，再加上$\lVert f\rVert_p$。$s$规定期望的变化阶数，$p$汇总同层位置，$q$汇总不同尺度；周期阶跃在$s<1/2$时属于每个$B^s_{2,q}$，在$s=1/2$时只属于$B^{1/2}_{2,\infty}$，清楚地分开了$s$与$q$的作用。完备性先在$L^p$中取极限，再借差分的连续性与Fatou引理把差分估计传到极限。小波刻画把同样的要求改写为系数条件：第$j$层的加权块大小$2^{j(s+1/2-1/p)}\lVert c_j\rVert_{\ell^p}$按$\ell^q$汇总，再加上粗项。等价性要求小波的正则性与消失矩都高于$s$，Haar小波只适用于较小的$s$。

在小波坐标中，三种逼近任务分别由不同条件控制。固定尺度截断为整类函数预先选定$2^J$维子空间，$B^s_{2,q}$球上的统一误差为$C_sRM^{-s}$，与$q$无关；反过来，截断误差序列$(2^{Js}E_J(f))_J$的$\ell^q$汇总恰好刻画$B^s_{2,q}$，单一的统一速率只对应$q=\infty$。最佳$N$项逼近让每个函数自选下标；在$1/\tau=s+1/2$时，$B^s_{\tau,\tau}$就是系数的$\ell^\tau$可和性，它给出$\sigma_N(f)_2\leq CN^{-s}$。误差速率本身只等价于弱$\ell^\tau$条件，不能反推Besov成员资格；把各个预算的误差按$\ell^\tau$汇总得到的逼近空间$A^s_\tau$，才与$B^s_{\tau,\tau}$一致。

整类函数的有限覆盖是第三种任务。对$B^s_{2,q}$球，先由统一截断界舍去尾部，再对保留的各层按几何方式分配精度，用有限维球的覆盖界逐层计数，得到$\log N(\varepsilon,\mathcal F_R)\leq C(R/\varepsilon)^{1/s}$；只在单一尺度上变化的子类是一个欧氏球，体积比较给出同阶的下界。与Sobolev紧嵌入只给出“有限”不同，这里得到了增长阶，它由$s$决定而与$q$无关。误差速率、误差汇总与对数覆盖数回答不同的问题，各自需要的参数条件不能相互替代。

最后，在单位周期上$B^1_{2,2}$与端点值相同的$H^1(0,1)$函数一致，范数等价；端点条件消去分部积分的边界项，二阶差分避开一阶连续模的饱和。$p\neq2$时整数阶Sobolev空间一般不与Besov空间重合。Besov尺度在Sobolev空间之外增加的，是分数的变化阶数、可选的尺度汇总，以及与稀疏逼近相匹配的积分指数。从离散观测出发、具有连续点值读取的函数空间，是第~\ref{chap:rkhs}章的对象。
